% Generated draft: requires mathematical and source review.
\documentclass[11pt]{article}
\usepackage[margin=25mm]{geometry}
\usepackage{fontspec}
\setmainfont{Latin Modern Roman}
\usepackage{amsmath,amssymb,mathtools,unicode-math}
\setmathfont{Latin Modern Math}
\usepackage{enumitem,needspace,xurl,hyperref}
\setlength{\parindent}{0pt}
\setlength{\parskip}{5pt}
\setlength{\emergencystretch}{3em}
\newcommand{\N}{\mathbb{N}}
\newcommand{\Z}{\mathbb{Z}}
\newcommand{\Q}{\mathbb{Q}}
\newcommand{\R}{\mathbb{R}}
\newcommand{\C}{\mathbb{C}}
\newcommand{\F}{\mathbb{F}}
\newcommand{\A}{\mathbb{A}}
\newcommand{\PP}{\mathbb{P}}
\newcommand{\E}{\mathbb{E}}
\newcommand{\eps}{\varepsilon}
\newcommand{\dd}{\,\mathrm{d}}
\begin{document}
\section*{UnsolvedMath ID: 5600007}
\textbf{Source problem code:} catalogue entry
\subsection{Definitions and mathematical statement}
Fix dimensions \(1\le n\le N\) and an area bound \(A>0\). The target \(S^N(1)\subseteq\R^{N+1}\) is the unit round sphere. For an immersion \(F:S^n\to S^N(1)\), let \(g_F=F^*g_{S^N}\). In the unrestricted submanifold interpretation, minimal means that the mean-curvature vector vanishes, and total area means the induced \(n\)-volume
\[
\operatorname{Vol}(F)=\int_{S^n}\mathrm d\operatorname{vol}_{g_F}\le A.
\]
Ambient congruence means \(G=Q\circ F\) for some \(Q\in O(N+1)\). Thus a literal candidate moduli set is
\[
\mathcal M_{n,N}(A)=
\{F:S^n\to S^N(1):F\text{ smooth minimal immersion},\ \operatorname{Vol}(F)\le A\}/O(N+1).
\]
Chern asks whether the intended moduli set is finite-dimensional with a natural topology. One candidate topology is the quotient of the smooth-map topology, but the source supplies neither that choice nor a dimension convention.

There is an essential metric/parametrization ambiguity. If no domain metric is fixed, precomposition by arbitrary diffeomorphisms of \(S^n\) preserves minimality and induced volume, and the stated ambient-motion equivalence does not remove that freedom. If instead the source means isometric minimal immersions of a fixed round sphere, these are a different class, with its radius/curvature needing specification. Quotienting by domain reparametrization is another possible formulation, but is not authorized by the stated equivalence. These versions are kept distinct rather than silently changing the moduli problem.

The supplied research summary labels the problem solved using bounded-degree compact convex parameter bodies and low-degree rigidity attributed to Calabi and do Carmo--Wallach. It does not state the domain-metric, isometry, degree, fullness, or equivalence hypotheses of those results, nor the implication from area bound to degree bound. Accordingly that assessment is preserved as supplied, but does not identify a precise theorem for the unrestricted candidate set above.
\subsection{Short English statement}
Are bounded-area minimal sphere immersions, modulo ambient motions, a finite-dimensional naturally topologized moduli space? The domain metric and treatment of reparametrization are essential to matching the reported literature solution.
\subsection{Sources}
\begin{itemize}
\item Local UnsolvedMath registry, permanent ID \(5600007\), exact source code \texttt{catalogue entry}, ``Moduli of Minimal Sphere Immersions.''
\item Shiing-Shen Chern, \emph{Open Problems in Differential Geometry} (1970), Problem 7, recorded PDF page 8; public ICM-proceedings scan at \url{https://people.math.harvard.edu/~hirolee/pdfs/2014-fall-230a-icm1970-chern-differential-geometry.pdf}. Recorded access 2026-07-29; pdf-text-normalized extraction.
\item Supplied research summary mentions Calabi and do Carmo--Wallach, compact convex parameter bodies, and low-degree rigidity without complete publication references or precise theorem hypotheses.
\end{itemize}
\end{document}
